\section{The bound ten under a six-cycle}
This section proves the six-by-six target when the whole level space is invariant under a full cycle of rows or columns. Invariance permits several character components. The proof controls their interactions through the quadratic form and the elementary-symmetric equations.

\begin{theorem}\label{thm:rowcycle}
Let $T$ cyclically permute the six rows of a complex $6\times6$ matrix. If $L$ is a level space and $T(L)=L$, then $\dim L\leq10$. The same statement holds for a full six-cycle of columns.
\end{theorem}

We first prove the small elementary-symmetric assertions on which the theorem rests.
\begin{lemma}\label{lem:singular-e3}
For $m\geq3$, the affine singular locus of $e_3$ in $\C^m$ is the union of the coordinate axes. The cubic hypersurface $e_3=0$ in $\mathbb P^4$ has five ordinary double points and contains no projective plane.
\end{lemma}
\begin{proof}
Put $s=\sum x_i$ and $e=e_2(x)$. The derivative in coordinate $i$ is $e-sx_i+x_i^2$. At a singular point, every coordinate is a root of one quadratic, so there are at most two values. If there is only one value, the derivative is $\binom{m-1}{2}$ times its square, and it is zero only at the origin. If the two values are $a,b$, occurring $r,m-r$ times, their sum and product are $s,e$. Comparing the first and second power sums gives
\[
 (r-1)a+(m-r-1)b=0,
 \qquad (r-1)a^2+(m-r-1)b^2=0.
\]
When both integer coefficients are positive, substitution forces $a=b=0$. If one coefficient is zero, the other value is zero and only one coordinate can be nonzero. Conversely every coordinate axis is singular.

For $m=5$, in the chart at a coordinate point the local equation is $e_2(z_1,\ldots,z_4)+e_3(z_1,\ldots,z_4)$. The quadratic part has matrix $(J_4-I_4)/2$, which is nondegenerate. Thus these five points are ordinary double points.

Suppose a projective plane $\Pi$ were contained in the cubic. In linear coordinates with $\Pi=\{y_1=y_2=0\}$, write the cubic as $y_1Q_1+y_2Q_2$. Its singular points on $\Pi$ are exactly the intersection of the two plane quadrics $q_i=Q_i|_\Pi$. This intersection is finite, since the entire singular locus is finite, so the quadrics form a regular sequence. Their Hilbert series is
\[
 \frac{(1-z^2)^2}{(1-z)^3}=\frac{(1+z)^2}{1-z},
\]
and their zero-dimensional projective intersection has length four. This Hilbert-series identity follows successively from the injective multiplication maps by the two quadrics and needs no genericity assumption.

Each of these intersection points is reduced. At an ordinary double point the four-dimensional tangent space has a nondegenerate quadratic form. The two-dimensional tangent space to $\Pi$ is totally isotropic. Its pairing with a complementary normal space is therefore nonsingular: a vector in the kernel would be orthogonal to both summands and hence zero. This is the Jacobian of $(q_1,q_2)$ along $\Pi$, so each local intersection has length one. It follows that $\Pi$ contains four distinct coordinate points. They span a projective three-space, a contradiction.
\end{proof}

\begin{lemma}\label{lem:evenelementary}
A complex linear subspace of $\C^5$ on which $e_2=e_4=0$ has dimension at most one. A subspace of $\C^6$ with $e_3=e_6=0$ has dimension at most two, and one with $e_2=e_4=e_6=0$ has dimension at most one.
\end{lemma}
\begin{proof}
For the first assertion regard the five coordinates as linear forms on the subspace. If one is identically zero, the equation $e_4=0$ forces a second to be zero, since a product of four nonzero linear forms cannot vanish identically. There remain at most three coordinates; the quadratic form $e_2$ on three coordinates is nondegenerate and has maximal complex isotropic dimension one. The cases of one or two remaining coordinates give the same upper bound directly.

If all five forms are nonzero, divide $e_4=0$ by their product in the fraction field. Grouping proportional forms as $a_iL_C$ gives
\[
 \sum_C\frac{\sum_{i\in C}a_i^{-1}}{L_C}=0.
\]
Multiply by the product of the distinct $L_C$ and reduce modulo each $L_C$. Pairwise relative primeness forces every numerator to be zero. Thus each proportionality class has at least two members, and there are at most two classes. The coordinate forms span the dual of the subspace, so its dimension is at most two. In dimension two the classes have sizes two and three. Write the forms as $aU,-aU,bV,cV,dV$. The reciprocal relation gives $bc+bd+cd=0$. In $e_2$, the cross term vanishes and the $V^2$ coefficient is zero, whereas the $U^2$ coefficient is $-a^2\ne0$. This contradiction proves the first assertion.

On a linear subspace of $\C^6$, the identity $e_6=0$ says that a product of coordinate linear forms is zero in an integral domain. Some coordinate is identically zero. The second assertion is then Lemma~\ref{lem:singular-e3} in the resulting five coordinates, and the third is the first assertion.
\end{proof}

\begin{proof}[Proof of Theorem~\ref{thm:rowcycle}]
Let $\zeta$ be a primitive sixth root of unity and let $v_b=(1,\zeta^b,\ldots,\zeta^{5b})^{\mathsf T}$. The spectral projections of $T$ preserve $L$, so
\[
 L=\bigoplus_{b=0}^5 v_b\otimes U_b
\]
for subspaces $U_b\subseteq\C^6$. Write $d_b=\dim U_b$. Every component is itself level. The rank-one formula of Proposition~\ref{prop:rankone} implies $U_0=0$ and $e_6|_{U_b}=0$ for every $b$. It also gives $e_3=0$ on $U_2,U_4$, and $e_2=e_4=0$ on $U_3$, using
\[
 \prod_{i=0}^5(1+z\zeta^{2i})=(1+z^3)^2,
 \quad\prod_{i=0}^5(1+z(-1)^i)=(1-z^2)^3.
\]
The two lemmas therefore give $d_2,d_4\leq2$ and $d_3\leq1$.

Polarizing $S_2$ between modes one and five gives
\[
 \beta(U_1,U_5)=0,\qquad
 \beta(u,w)=\left(\sum_j u_j\right)\left(\sum_jw_j\right)-\sum_ju_jw_j.
\]
For example the cross coefficient is $-6\beta(u,w)$, obtained by summing $\zeta^i\zeta^{-i'}$ over distinct row indices. The matrix $B=J_6-I_6$ is invertible with $B^{-1}=-I_6+J_6/5$, so $d_1+d_5\leq6$. Equality would imply $U_5=U_1^{\perp_\beta}$. As $e_6|_{U_1}=0$, some coordinate $j$ vanishes throughout $U_1$. Then $B^{-1}e_j\in U_5$, but its entries are $-4/5$ at $j$ and $1/5$ elsewhere; its sixth elementary symmetric function is nonzero. This contradicts $e_6|_{U_5}=0$. Hence $d_1+d_5\leq5$, and
\[
 \dim L\leq5+2+1+2=10.
\]
Transposition proves the column statement.
\end{proof}

\begin{corollary}\label{cor:defect}
For any $6\times6$ level space $L$ and full row six-cycle $T$, put
$r=\dim L-\dim(L\cap TL)$. Then $\dim L\leq10+5r$.
\end{corollary}
\begin{proof}
Intersect $L,TL,\ldots,T^5L$ successively. At each step the dimension drops by at most $r$: the preceding intersection lies in the preceding translate, whose intersection with the next translate has codimension $r$. Their full intersection $K$ has dimension at least $\dim L-5r$, is $T$-invariant, and is level. Apply Theorem~\ref{thm:rowcycle}.
\end{proof}

\begin{proposition}\label{prop:twotransitive}
If a finite two-transitive group on $n\geq2$ row positions preserves a level space $L$, then $\dim L\leq n-1$. The bound is attained by the sum-zero matrices supported in one column.
\end{proposition}
\begin{proof}
The permutation module is $\C\mathbf1\oplus V$, with $V$ the irreducible sum-zero module. Indeed its commuting algebra has dimension two, the number of orbits on ordered pairs, and complete reducibility makes this the sum of the squares of irreducible multiplicities. Therefore any invariant subspace of its tensor product with the column space has the form
$\C\mathbf1\otimes U_0\oplus V\otimes U_1$.
The rank-one conditions force $U_0=0$. For every $2\leq k\leq n$ one can choose a vector of $V$ with nonzero $e_k$: use all $k$th roots of unity in $k$ coordinates and zeros elsewhere. Thus $e_k|_{U_1}=0$ for every $k\geq2$. The polynomial with the coordinates of $u\in U_1$ as roots is $z^n-e_1(u)z^{n-1}$, so $u$ has at most one nonzero coordinate. A linear subspace in a union of coordinate axes lies in one axis, giving $\dim U_1\leq1$ and the result.
\end{proof}

These invariant results do not arise by averaging an arbitrary space. The row-supported level direction $(1,-1,0,\ldots,0)$ has cyclic average with every row $(1,-1,0,\ldots,0)/n$; its second rook number is $-(n-1)/n$, so the average is not level.

\section{Scalar symmetries and maximal hooks}
This section excludes two other ways of producing an eleven-dimensional example.

\begin{theorem}\label{thm:scalar-six}
Let $L$ be a $6\times6$ level space. If some $g\in(S_6\times S_6)\rtimes\langle\mathsf T\rangle$ acts as a scalar $\omega\ne1$ on $L$, then $\dim L\leq10$.
\end{theorem}
\begin{proof}
Let $q$ be the order of $\omega$. For $q\geq4$ the orbit count gives ambient eigenspace dimension at most $\lfloor36/q\rfloor\leq9$. If $q=2$, the $-1$ eigenspace lies in $S_1=0$. Since $g$ preserves its nondegenerate quadratic form $S_2$, that eigenspace is orthogonal to every other eigenspace and its restricted form is nondegenerate. Its dimension is at most eighteen, so its totally isotropic subspace $L$ has dimension at most nine.

Consider $q=3$ without transposition. A row cycle of length $a$ and a column cycle of length $b$ give $\gcd(a,b)$ coordinate cycles of length $\operatorname{lcm}(a,b)$. A primitive cube root therefore has multiplicity
\[
 \sum_{3\mid\operatorname{lcm}(a,b)}\gcd(a,b).
\]
Substitution of the eleven partitions of six shows that multiplicity at least eleven must equal twelve. After exchanging rows and columns, the row permutation has type $3+3$ and the column type is $1^6$, $3+1^3$, or $3+3$. To make this finite check explicit, the maximum multiplicities over all column types for row types
\[
 6,\ 5+1,\ 4+2,\ 4+1+1,\ 3+3,\ 3+2+1,
 \ 3+1+1+1,\ 2+2+2,\ 2+2+1+1,\ 2+1^4,\ 1^6
\]
are respectively $6,4,4,6,12,10,12,6,8,10,12$; the formula also identifies the pairs giving twelve.

In each twelve-dimensional case the restricted cubic $S_3$ has no linear factor. For identity columns, restrict to one row triple and four free columns. The resulting polynomial is $6e_3(a,b,c,d)$. The polynomial $e_3=d(ab+ac+bc)+abc$ is primitive and linear in $d$, with relatively prime coefficients in $\C[a,b,c]$, so it is irreducible. If the column permutation contains a three-cycle, restrict to one row triple and one column triple. The matrix is a circulant $3\times3$ matrix with rows multiplied by $1,\omega,\omega^2$, and its permanent is
\[
 a^3+b^3+c^3+3abc.
\]
This plane cubic is smooth: its gradient equations give $a^2=-bc$, $b^2=-ac$, $c^2=-ab$; multiplying rules out a nonzero solution, and a zero coordinate forces all three to vanish. A reducible plane cubic has intersecting components and is singular, so this polynomial has no linear factor either. A linear factor of the full restricted cubic would restrict to zero, contradicting the nonzero restricted polynomial, or to a linear factor of it. An eleven-dimensional subspace in a twelve-dimensional space is a hyperplane; vanishing there would force just such a factor.

For an element involving transposition, conjugate its coordinate action to $(i,j)\mapsto(j,\pi(i))$. Between distinct cycles of $\pi$ of lengths $a,b$, the two rectangles give $\gcd(a,b)$ cycles of length $2\operatorname{lcm}(a,b)$. Within a cycle square there are $\lfloor a/2\rfloor$ cycles of length $2a$, and one additional cycle of length $a$ if $a$ is odd. Label the cycle by $\Z/a$ to check this: squaring the map translates both coordinates, while their difference is sent to one minus itself. For the eleven cycle types above, the primitive-cube-root multiplicities are
$3,0,0,0,7,4,5,0,0,0,0$. They are at most seven. This completes every case.
\end{proof}

\begin{theorem}\label{thm:hook}
For $p,q\geq2$, let $H_{p,q}$ be the matrices supported in the first row and first column whose first row sums to zero and whose first column below that row sums to zero. If a linear subspace $L$ contains $H_{p,q}$ and $S_1,S_2,S_3$ vanish on $L$ (omitting orders above $\min(p,q)$), then $L=H_{p,q}$. In particular this level space of dimension $p+q-3$ is maximal by inclusion.
\end{theorem}
\begin{proof}
The matching-support argument already proves that $H_{p,q}$ is level. Take $A\in L$ and vary a sum-zero first row $r$. The coefficient of $t$ in $S_2(A+tr)$ is
\[
 \sum_jr_{1j}\sum_{i>1,\ell\ne j}A_{i\ell}.
\]
It follows that $v_j=\sum_{i>1}A_{ij}$ is independent of $j$. Varying a sum-zero column $c$ below the first row likewise makes $R_i=\sum_{j>1}A_{ij}$ independent of $i>1$.
Let $B$ be the block outside the first row and column. When $p,q\geq3$, the coefficient of $tu$ in $S_3(A+tr+uc)$ is
\[
 \sum_{i>1,j>1}c_{i1}r_{1j}T_{ij},
 \quad T_{ij}=\sum_{a>1,\ a\ne i,\ b>1,\ b\ne j}B_{ab}.
\]
The entries $r_{1j}$ for $j>1$ are free, while $c$ ranges through its sum-zero space. Hence $T_{ij}$ is independent of $i$ for each $j$. But
$T_{ij}=\sum B-R_i-v_j+B_{ij}$, so $B_{ij}$ is independent of $i$; equality of its column sums makes every entry one scalar $b$. If $p=2$, this follows from equality of the $v_j$, and if $q=2$ from equality of the $R_i$.

The condition $S_1(A)=0$ now shows that
$A-b(J_{p,q}-pqE_{11})\in H_{p,q}$. Thus $b(J_{p,q}-pqE_{11})\in L$. Counting two-edge matchings gives
\[
 S_2\bigl(b(J_{p,q}-pqE_{11})\bigr)
 =-\tfrac12pq(p-1)(q-1)b^2.
\]
Its coefficient is nonzero, so $b=0$ and $A\in H_{p,q}$.
\end{proof}

Maximality by inclusion does not assert maximum dimension. In particular the theorem rules out extending the standard nine-dimensional hook to dimension eleven, but does not classify all eleven-dimensional subspaces.
